% Titelseite der eigenen Anlageanalyse. Bewusst ohne Hochschulwappen,
% Matrikelnummer und Studiengangsangabe: Das Dokument ist keine Pruefungs-
% leistung, sondern eine Entscheidungsgrundlage fuer eigene Anlagen.
% Der Bearbeitungsstand steht auf der Titelseite, weil eine Analyse mit
% jedem neuen Quartalsbericht altert.
\begin{titlepage}
\centering
\vspace*{\stretch{2}}
{\large Anlageanalyse aus Prim\"arquellen\par}
\vspace{1.2cm}
{\Huge\bfseries Alamos Gold Inc.\par}
\vspace{0.5cm}
{\LARGE NYSE und TSX: AGI\par}
\vspace{0.4cm}
{\Large Gesch\"aftsjahr 2025\par}
\vspace{\stretch{2}}
{\large Unternehmen, Jahresabschluss, strategische Ma{\ss}nahmen,\\
Kursverlauf, Flow-Analyse und Anlageurteil\par}
\vspace{\stretch{3}}
\begin{tabular}{ll}
Verfasser: & Yann Bertin Kamdem Bobda\\
Stand:     & 16.~September 2026\\
Quellen:   & Form 40-F 2016--2025, Zwischenbericht zum 30.06.2026,\\
           & Management Information Circular 2026\\
\end{tabular}
% \par schliesst den Absatz der Tabelle ab. Ohne ihn setzt LaTeX den
% folgenden Hinweis in dieselbe Zeile wie die letzte Tabellenzeile, und
% \vspace wirkt dann als horizontaler statt vertikaler Abstand.
\par
\vspace{1.2cm}
\begin{minipage}{0.8\linewidth}
\centering\small
Alle Zahlenangaben aus Prim\"arquellen sind gegen die gedruckte Seite ihrer
Quelle gepr\"uft. Blau gesetzte Passagen sind Einsch\"atzungen, rot gesetzte
sind das Votum in Teil~\ref{sec:verdikt}. Beide sind keine Belege.
\end{minipage}
\vspace*{\stretch{2}}
\end{titlepage}
